\section{Introduction}
Patients are eligible to be included in the datset

\subsection{Hypothetical trial protocol}
\paragraph{Eligibility criteria} Veterans are eligible to enroll in the trial if they meet each of the following criteria:
\begin{enumerate}
    \item Was a Veteran who was $\geq$ 18 years old
    \item Contacted the National Call Center for Homeless Veterans (NCCHV) for the first time between XX/XXXX and XX/XXXX
    \item Was deemed to be either homeless or at risk of homelessness at the time of the call.
    \item Had never enrolled in SSVF prior to their Call Center contact.
    \item Was not enrolled in any VA homeless program at the time of their initial Call Center contact.
\end{enumerate}

\paragraph{Treatment strategies.} Upon trial enrollment, participants would be randomly assigned to one of two arms: \textit{SSVF} or \textit{No SSVF}. Participants in the \textit{SSVF} arm would be instructed to enroll in SSVF within 30 days of their Call Center contact (the grace period). We define $\Delta_A = \mathbf{1}[T_A \leq 30]$ as the indicator that a participant initiated SSVF within the grace period.

\paragraph{Follow-up.} Participants would be followed for 120 days after treatment initiation. The primary outcome is a composite indicator of being alive and housed at 120 days after SSVF initiation: $Y(T_A + 120) = (1 - D(T_A + 120)) \cdot H(T_A + 120)$.
